\begingroup\sloppy
\section{Questions ouvertes et suite}\label{sec:ouvertes-ref}\label{ouv:sec}

Cet inventaire rassemble les questions ouvertes du ledger \texttt{v2/\allowbreak{}docs/\allowbreak{}lab/\allowbreak{}open.md}, de la section \enquote{Open / not yet decided} de \texttt{v2/\allowbreak{}docs/\allowbreak{}lab/\allowbreak{}decisions.md}, et des actions proposées par les rapports du 2026-11-15 (audit de conformité, grille du jury) qui ne figurent pas encore dans le ledger.
Les doublons sont fusionnés : une ligne du tableau peut regrouper plusieurs lignes du ledger.

\paragraph{Lecture des colonnes.}
\emph{Prix} : coût estimé par le ledger ou par le rapport source, en temps Claude sauf mention contraire ; \enquote{humain} désigne le temps d'une personne de l'équipe ou d'un tiers (DPO, juriste, ARS), \enquote{monde} un délai qui ne dépend pas de l'équipe (publication d'un millésime, réponse d'un producteur). Tous les prix sont des estimations\suppose{}.
\emph{Priorité} : jugement du présent rapport. \textbf{H} : bloque la reprise, une affirmation publique ou un chiffre cité ; \textbf{M} : améliore la précision ou le produit ; \textbf{B} : confort ou nettoyage.
\emph{Source} : fichier où la question est posée.

\paragraph{Questions fermées.}
Les questions tranchées ne figurent pas dans le tableau\footnote{Sur environ 135\derive{} lignes d'inventaire (après une première fusion des doublons), 32\derive{} ont été écartées comme fermées, et 2\derive{} autres comme fermées en partie (seule la partie ouverte est reprise). Une ligne est fermée quand \texttt{open.md} le dit, ou quand une entrée datée de \texttt{decisions.md} la tranche. Environ une douzaine\derive{} de ces dernières sont encore marquées \enquote{open} dans \texttt{open.md} : c'est une incohérence du ledger (annexe~\ref{pro:inc-e}). Le classement a été fait à la main, à la lecture du registre ; un auteur d'axe peut juger une question encore ouverte.}.

\begin{small}
\setlength{\LTleft}{0pt}
\begin{longtable}{@{}p{0.48\textwidth}p{0.16\textwidth}p{0.05\textwidth}p{0.21\textwidth}@{}}
\caption{Questions ouvertes, dédoublonnées et groupées par thème.}\label{ouv:tab}\\
\toprule Question & Prix & Prio. & Source \\ \midrule\endfirsthead
\toprule Question & Prix & Prio. & Source \\ \midrule\endhead
\bottomrule\endfoot
% ------------------------------------------------------------------
\multicolumn{4}{@{}l}{\textbf{Données}}\\
Part privée des pharmaciens hospitaliers : RPPS 16,6--17,4\mesure{}\,\%, SAE 29,4\mesure{}\,\%, Ordre 28,6\mesure{}\,\% ; premier facteur d'incertitude des pharmaciens à l'axe 5 ($S_T=0{,}48$\mesure{}) & 30\suppose{}~min à 1\suppose{}~h & H & \texttt{open.md} (3 lignes), \texttt{decisions.md} § Open \\
Brancher \texttt{secteur\_\allowbreak{}stock.parquet} dans l'état du jumeau & 30\suppose{}~min & M & \texttt{open.md} \\
Ventiler \enquote{autre secteur} : pharmaciens salariés (5\,805\mesure{}, part 0,5\suppose{}) et sages-femmes (796\mesure{}, part 0,5\suppose{}) & 1\suppose{}~h chacun & M & \texttt{open.md} \\
Parts privées FINESS pondérées par effectifs (centres de santé, médico-social) & 1\suppose{}~h~30 & M & \texttt{open.md} \\
Classer \enquote{Prévention} en public pour les sages-femmes (écart min--max $-51$\derive{}\,\%) & 15\suppose{}~min & M & \texttt{open.md} \\
Double compte des infirmiers libéraux (SNDS) et salariés (BTS) ; hypothèse $H_\mathrm{ind}$ non testable sur l'âge & 30\suppose{} + 45\suppose{}~min & M & \texttt{open.md} \\
Dater les zonages : Atlasanté \texttt{c987d0e2} (2012 ou 2014), CartoSanté actuel par ARS (17\mesure{}), zonages intermédiaires 2018--2023 et de 1\textsuperscript{re} génération & 30\suppose{}~min à 1\suppose{}~h chacun, issue incertaine & M & \texttt{open.md} \\
Textes primaires des dates d'effet (avenants IDE, rezonage, dentaire, avenant~7 des kinés par ARS) et des 12\mesure{} événements à source secondaire & 30\suppose{}~min à 1\suppose{}~h chacun & M & \texttt{open.md} \\
Bornes des objectifs \emph{numerus apertus} 2021--2025, bilan ONDPS & 30\suppose{}~min & M & \texttt{open.md} \\
Inscrits en formation : artefact Paris Cité 2025-26 (683\mesure{} au lieu d'environ 190\suppose{}) ; série par université avant 2020 & 15\suppose{}~min + monde (DREES) ; 1\suppose{}~h & M & \texttt{open.md} \\
Écart CartoSanté contre RPPS 2025 (kinés $-5{,}4$\mesure{}\,\%, Marseille absente) & 30\suppose{}~min & M & \texttt{open.md} \\
Archiver l'Annuaire Santé (au moins 3\suppose{} instantanés ; le premier date du 2026-11-13) & 10 à 30\suppose{}~min + monde ; décision & H & \texttt{open.md} (3 lignes) \\
Historique communal : proxy d'installations nettes des moins de 40 ans, flux d'installation et de cessation, pharmacies 2004--2025 (FINESS historique), sages-femmes (aucune histoire communale ouverte) & 45\suppose{}~min à 1\suppose{}~h chacun ; SF : monde & M & \texttt{open.md} \\
Table de voisinage communale (destination \enquote{zone voisine}) ; table de passage RPPS privé $\leftrightarrow$ CartoSanté libéral & 30\suppose{}~min ; 1\suppose{}~h & M & \texttt{open.md} \\
Diplômés hors de France : flux annuel par profession depuis une source primaire (défauts proposés, non mesurés : 1\,300 contre 665\suppose{} pour les dentistes) & 45\suppose{}~min à 1\suppose{}~h~30 & H & \texttt{open.md} (4 lignes), \texttt{decisions.md} \\
Vérifier le grain le plus fin d'Omphale et d'Hippocrate & 20\suppose{}~min & B & \texttt{open.md} \\
Taux d'installation au-dessus de 100\,\% (kinés 107\mesure{}\,\%) : renommer, ne compter que les premières entrées, ou masquer & 10\suppose{}~min à 2\suppose{}~h & M & \texttt{decisions.md} § Open \\
% ------------------------------------------------------------------
\multicolumn{4}{@{}l}{\textbf{Modèle}}\\
Tester la règle \enquote{entrées $\leq$ sorties} contre une référence indépendante (C13 ne la teste qu'en interne) & 30\suppose{}~min & H & \texttt{open.md} \\
Vérifier 3\mesure{} constats non réfutés de la revue (conservation, vieillissement d'une tranche sur 5, état des modules) & 10\suppose{}~min & M & \texttt{open.md} \\
Séparer entrées et rétention (abandonner \enquote{excédent de $r>1$ = entrée}) ; cause commune de M3, M4 et de l'échec SBC & 2\suppose{}~h & H & \texttt{open.md} \\
Sensibilité de $N(2050)$ à l'ordre des modules ; plancher à 0 hérité de la v1 (0,02 à 0,43\mesure{}\,\% de $N(2050)$) ; réaligner \texttt{decalage\_\allowbreak{}retraite} & 10 à 20\suppose{}~min chacun & B & \texttt{open.md}, \texttt{decisions.md} § Open \\
État communal projetable (tranches CartoSanté variables, \texttt{fit\_\allowbreak{}rates} en 1/5) ; recalcul communal 5 professions $\times$ 2 secteurs (1,1\mesure{}~s mesurée, 3,0\derive{}~s extrapolée) & 20\suppose{}~min à 1\suppose{}~h & M & \texttt{open.md} \\
Injecter les diplômés du vivier dans $N$ à la place du flux ajusté ; variante \texttt{vivier} au banc N1 ; abandon en formation IDE projeté ($\rho$ de 0,84 à 0,72\mesure{}) & 30\suppose{}~min à 1\suppose{}~h ; 45\suppose{}~min & M & \texttt{open.md} \\
Âge d'entrée des diplômés à l'étranger : 28\suppose{} ans unique, ou un âge par profession (effet $\times2$ : dentistes $-1{,}3$\mesure{}\,\%, pharmaciens $-0{,}2$\mesure{}\,\%) & 30\suppose{}~min & B & \texttt{decisions.md} § Open \\
Contrat du jumeau : \texttt{Shift} sans fraction par cellule, \texttt{stock\_\allowbreak{}flow\_\allowbreak{}rates} non injectable, point d'injection officiel des taux de C1 & 20 à 30\suppose{}~min chacun & M & \texttt{open.md} \\
Reprise de l'observatoire vers l'axe 0 (B1, B2, B6, A6) & 1\suppose{}~h~30 & B & \texttt{open.md} \\
Axe 1 : paramètres à sourcer (arrêtés de zonage, avenant~11 IDE, taux de dérogation, $\varepsilon_\mathrm{PH}$, postes de DES d'odontologie, cessation par âge, option \enquote{2 ans} des kinés non mesurable avant 2028) ; relecture (\texttt{param-\allowbreak{}regul-\allowbreak{}report-\allowbreak{}entrants}, DREES Écoles 1E, ONCD 1\,333 ou 1\,321\mesure{}) & 10\suppose{}~min à 1\suppose{}~h chacun & M & \texttt{open.md} \\
Effet chiffré de la régulation pour SF, MK et CD ; destination des entrants bloqués ; régulation SF dérogatoire ; part des dentistes spécialistes exemptés & 20\suppose{}~min à 4\suppose{}~h & H & \texttt{open.md} \\
Rétention dépendante de la densité pour les IDE (N9 : $+0{,}11$\mesure{} [$+0{,}05$\,;\,$+0{,}18$], corrélationnel) & non chiffré & M & \texttt{decisions.md} § Open \\
Unité de l'offre : effectif ou ETP ; groupes de communes (TDS, zonage ARS, bassins de vie, EPCI) & non chiffré ; décision & M & \texttt{decisions.md} § Open (hérité v1) \\
Axe 2, bancs : A2 + A1b au banc N1 ; IC de $\Delta(\mathrm{A2}-\mathrm{A0})$ par blocs ; parts privées min et max ; relancer N1 et N6 sur la table privée actuelle (SF : A2 8,60 au lieu de 8,57\mesure{}) & 5 à 20\suppose{}~min chacun & M & \texttt{open.md} \\
Axe 2, C1 : entrées probabilistes ; bayésien complet de $\sigma_u,\sigma_v,\tau_r$ ; gain décroissant sur les origines récentes ; rétention variable avec l'horizon ; C1 sur la table privée IDE & 20\suppose{}~min à 1\suppose{}~j & M & \texttt{open.md}, \texttt{decisions.md} \\
Axe 2, autres familles : logit d'installation A4, espace-état A5 à la commune, multi-états privé $\leftrightarrow$ public ; GBM direct à $h=5$ ou sur les résidus de A2 (le GBM global est rejeté) & 30\suppose{}~min à 3\suppose{}~h & B & \texttt{open.md}, \texttt{2026-\allowbreak{}11-\allowbreak{}15-\allowbreak{}c1-\allowbreak{}gbm-\allowbreak{}global.md} \\
Budget d'erreur : covariables de l'écart local ($R^2\geq0{,}2$\suppose{} exigé), variance sur-dispersée (MSE $\approx6\times$\mesure{} poissonien), plancher de bruit sur flux bruts, persistance sur plus de paires & 10\suppose{}~min à 1\suppose{}~h~30 & M & \texttt{2026-\allowbreak{}11-\allowbreak{}15-\allowbreak{}budget-\allowbreak{}erreur-\allowbreak{}c1-\allowbreak{}a2.md} \\
Pharmaciens communaux : erreur médiane de 10,7\mesure{}\,\% à l'EPCI pour un seuil de 10\suppose{}\,\% & 1\suppose{}~h & B & \texttt{open.md} \\
Vectoriser le vivier (60\mesure{}\,\% du temps d'un run) & 1\suppose{}~h & B & \texttt{open.md} \\
Écarts aux références DREES : IFSI, perte divisée par 2 ($+15{,}6$\mesure{}\,\% contre $\approx+7{,}9$\,\% DREES) ; retraite $+1$ an à relancer sur le champ DREES & 45\suppose{}~min ; 30\suppose{}~min & M & \texttt{open.md} \\
Kinés quotas $+20$\,\% en 2050 : $+19\,110$\mesure{} aujourd'hui contre $+8\,645$\mesure{} dans le rapport du prototype ; mettre le rapport à jour, en rappelant qu'aucun chiffre 2050 n'est citable (\enquote{Axe 3, nuit}) & 20\suppose{}~min & H & \texttt{open.md} \\
% ------------------------------------------------------------------
\multicolumn{4}{@{}l}{\textbf{Validation}}\\
N2-A fidèle : quotas 2016--2020 ouverts ; tolérance de $\pm3$\suppose{}\,\% à 2050, bien plus étroite que les fourchettes de l'axe 5 & 45\suppose{}~min ; 15\suppose{}~min + décision & H & \texttt{open.md} \\
A0 et B4 divergent de $+23$\mesure{}\,\% pour les dentistes à N2 & 30\suppose{}~min & M & \texttt{open.md} \\
Relancer le banc K1 + filière sur le périmètre \enquote{libéral} (privé) & 20\suppose{}~min & B & \texttt{open.md} \\
N3 : 3\textsuperscript{e} test de la régulation IDE par contrôle synthétique ; pré-enregistrer $\theta$ pour E2 kinés & 1\suppose{}~h ; 15\suppose{}~min & H & \texttt{open.md}, \texttt{decisions.md} (\enquote{Axe 3, nuit}) \\
N4 : projection réelle à l'échelle fine ; borne $r\leq1{,}2$\suppose{} des mondes simulés & 15\suppose{}~min ; non chiffré & M & \texttt{open.md} \\
Axe 5 : composante (b) trop large à $h=5$ ; doubler $N$ (512 $\rightarrow$ 1\,024\suppose{} ; indice PH 0,085\mesure{} $>0{,}05$) ; Sobol par département ; $H_\mathrm{ind}$ comme facteur ; relance à chaque changement du jumeau & 20\suppose{}~min à 2\suppose{}~h machine & M & \texttt{open.md} \\
C1 : choc commun national ; intervalles sur une 2\textsuperscript{e} origine (2027) ; conforme par taille ; emballage en \texttt{Estimator} & 10\suppose{}~min à 2\suppose{}~h & M & \texttt{open.md}, \texttt{2026-\allowbreak{}11-\allowbreak{}15-\allowbreak{}c1-\allowbreak{}intervalles.md} \\
Scorer les prévisions déposées 2027--2030 (empreinte \texttt{ca036374d216}) & monde (juillet 2027) & H & \texttt{decisions.md} \\
Vérifier Gneiting et Ranjan (2011) en texte intégral & 15\suppose{}~min & B & \texttt{open.md} \\
N9 B2 avec flux bruts d'installation & 1\suppose{}~h & B & \texttt{open.md} \\
Réconcilier les IDE 2050 avec ER1319 (821\,000\mesure{})~\citep{e-er1319} et le \enquote{libéral} avec le RPPS (condition n\textsuperscript{o}~4 du rapport produit) & 20\suppose{}~min à 1\suppose{}~h & H & \texttt{open.md} \\
% ------------------------------------------------------------------
\multicolumn{4}{@{}l}{\textbf{Produit}}\\
\textbf{LLM ou déterministe} : utilité comparée sur les 22\mesure{} questions, règle d'ancrage stricte (le mode IA est déjà par défaut, §\ref{pro:llm-det}) & 30\suppose{}~min + 10\suppose{}~min machine + 3\suppose{}~\$ & H & \texttt{open.md}, \texttt{decisions.md} \\
Assistant : interdire une absence d'effet non calculée (cas \texttt{q06} et \texttt{h06}, §\ref{pro:cas-ecart}) & non chiffré & H & présent rapport ; \texttt{v2/\allowbreak{}agent/\allowbreak{}reponses\_\allowbreak{}*.jsonl} \\
Mesurer le mode IA sur les jeux 5 et 6 (compréhension, 60\mesure{} appels) et sur le périmètre (27\mesure{} cas) & 60 + 27\mesure{} appels & H & \texttt{chatbot/\allowbreak{}DECISIONS.md}, \texttt{open.md} \\
Resserrer la tolérance de $\pm6$\mesure{}\,\% du chatbot & 20\suppose{}~min & M & \texttt{open.md} \\
Jeux de test écrits par un humain : compilateur (et jeu 7 à l'aveugle), chatbot (avec relecture qualitative) & 30\suppose{}~min humain chacun ; 30\suppose{}~min & H & \texttt{open.md} \\
Glossaire : références juridiques à relire par un·e juriste & 30\suppose{}~min humain & M & \texttt{open.md} \\
Mistral imposé par le hackathon ? & décision & M & \texttt{decisions.md} § Open, \texttt{chatbot/\allowbreak{}DECISIONS.md} \\
Rejouer la recette par personas sur la carte finale (DSFR, option 1), avec re-recette ciblée des P0 et P1 corrigés et captures sur disque & 30 à 40\suppose{}~min & H & \texttt{open.md} \\
Recette avec des utilisateurs réels (ARS, administration, élu·e) & humain, non chiffré & H & présent rapport (aucune recette humaine, §\ref{pro:recettes}) \\
Accueil : entrée territoriale ou nationale, modes Simple et Expert & 30 à 45\suppose{}~min & M & \texttt{decisions.md} § Open \\
Latence perçue au premier plan, seuil \enquote{geste continu ou changement isolé} ; délai d'effet affiché 2032 ou 2033 ; exigence \enquote{très fluide} ($<100$\suppose{}~ms en rafale non établi) & 5 à 10\suppose{}~min + décision & B & \texttt{open.md} \\
Retour utilisateur sur la carte (zoom, cases, limites ARS, DROM) ; refonte du panneau Leviers ; direction artistique modifiable dans Figma & non chiffré & B & \texttt{open.md}, \texttt{decisions.md} § Open \\
Articulation avec l'observatoire de l'équipe (échange JSON), choix d'une direction artistique & 10\suppose{}~min humain & B & \texttt{open.md} \\
Corriger le deck ($\Delta$ MAPE, statut des leviers, diapo \enquote{Decisions needed}) et le \texttt{README.md} (annexe~\ref{pro:inc-e}) & 0,05 à 0,3\suppose{}~h & H & \texttt{2026-\allowbreak{}11-\allowbreak{}15-\allowbreak{}jury-\allowbreak{}grille.md} \\
Modèle économique (EUPL + prestations $<60$\suppose{}~k\texteuro{}~HT) ; porteur institutionnel et ARS pilote & 10\suppose{}~min humain ; semaines, humain & H & \texttt{2026-\allowbreak{}11-\allowbreak{}14-\allowbreak{}produit.md}, \texttt{open.md} \\
Relire en entier les 6\mesure{} pages sources du rapport produit & 20\suppose{}~min & B & \texttt{open.md} \\
Gabarit des prétirages : guillemets sous tectonic & 10\suppose{}~min & B & \texttt{open.md} \\
% ------------------------------------------------------------------
\multicolumn{4}{@{}l}{\textbf{Conformité}}\\
Seuil de secret statistique sur les sorties et les tables communales (2\mesure{} copies) : masquage, ou retrait du dépôt public & 30\suppose{}~min ou 20\suppose{}~min ; décision & H & \texttt{open.md}, audit \\
Note de non-applicabilité de l'AIPD & 1\suppose{}~h + DPO, humain & H & \texttt{open.md}, audit \\
Licence du code (EUPL~1.2 proposée) et \texttt{NOTICE} des licences tierces & 10\suppose{}~min humain + 15\suppose{}~min & H & \texttt{open.md}, audit \\
DSFR : nommer la structure porteuse (question levée le 2026-11-15 : projet porté par le Ministère chargé de la santé, selon l'utilisateur·rice) & 10\suppose{}~min humain & M & \texttt{decisions.md} § Open \\
Page \enquote{Conformité} du \texttt{README.md} (finalité, base légale, art.~14, destination non clinique, mention IA) & 20\suppose{}~min & H & audit (absent de \texttt{open.md}) \\
Avertissement IA (art.~50) et \enquote{pas de données personnelles} à l'ouverture de l'assistant & 10\suppose{}~min & H & audit (absent de \texttt{open.md}) \\
Conditions de réutilisation de CartoSanté ; licences des autres sources, fichier par fichier & 15\suppose{}~min ; non chiffré & M & audit (absent de \texttt{open.md}) \\
Relecture par un DPO ou un juriste & humain & H & audit \\
% ------------------------------------------------------------------
\multicolumn{4}{@{}l}{\textbf{Module de besoin} (détail en §\ref{ouv:besoin})}\\
Construire le module de besoin, sous réserve du test d'abandon & 1\suppose{}~h~30 (1 à 2\suppose{}~h~30), non calibré & H & \texttt{derivations/\allowbreak{}2026-\allowbreak{}11-\allowbreak{}14-\allowbreak{}module-\allowbreak{}besoin.md}, \texttt{decisions.md} § Open \\
Test d'abandon : besoin pondéré par l'âge contre besoin proportionnel aux habitants en 2050 (IDE, MK) & 30\suppose{}~min & H & idem \\
Sourcer $a$ et $k$ pour au moins un des trois transferts (L23, L25, L27) & non chiffré & H & idem \\
Recours observé contraint par l'offre : choisir une correction (M2) ou l'assumer & non chiffré ; décision & M & présent rapport \\
Reproduire l'APL DREES comme cible externe (M4) & non chiffré ; données d'activité manquantes & M & \texttt{archive-\allowbreak{}v1/\allowbreak{}explorations/\allowbreak{}drees-\allowbreak{}projections-\allowbreak{}and-\allowbreak{}apl.md} \\
Médecine générale comme \enquote{source} d'un transfert : dans ou hors du périmètre ? & 10\suppose{}~min humain ; décision & M & présent rapport \\
% ------------------------------------------------------------------
\multicolumn{4}{@{}l}{\textbf{Spécialités : amélioration possible (hors plan depuis le 2026-11-15)}}\\
Cardiologie, dermatologie, ophtalmologie : postes d'internat par DES, secteur 1 ou 2, dépassements. Condition de démarrage (axes 0 à 3 terminés pour au moins une profession, dont un levier validé) non remplie & non chiffré & M & \texttt{v2/\allowbreak{}ROADMAP.md}, \texttt{open.md} \\
Spécialistes comme \enquote{source} d'un transfert de tâches (contexte du module de besoin) & inclus dans le module & B & \texttt{derivations/\allowbreak{}2026-\allowbreak{}11-\allowbreak{}14-\allowbreak{}module-\allowbreak{}besoin.md} \\
\end{longtable}
\end{small}

% ---------------------------------------------------------------------------
\subsection{Module de besoin}\label{ouv:besoin}

\paragraph{Statut.}
Plan écrit, non construit, en attente d'accord (\texttt{v2/\allowbreak{}docs/\allowbreak{}lab/\allowbreak{}derivations/\allowbreak{}2026-\allowbreak{}11-\allowbreak{}14-\allowbreak{}module-\allowbreak{}besoin.md}). Le dossier \texttt{v2/\allowbreak{}besoin/\allowbreak{}} n'existe pas.
La question est au registre (\texttt{decisions.md}, § Open, \enquote{Leviers de demande}).
Le jury estime qu'un indicateur de besoin rapporterait environ $+2{,}5$\suppose{} points pour 3\suppose{}~h, et sa question dure n\textsuperscript{o}~15 (\enquote{où est la demande ?}) reste sans réponse (\texttt{2026-\allowbreak{}11-\allowbreak{}15-\allowbreak{}jury-\allowbreak{}grille.md}).

\paragraph{But.}
Rendre répondable la question du sujet, \enquote{combien de spécialistes à Clermont-Ferrand en 2035 si les pharmaciens peuvent vacciner ?}, \emph{sans changer la projection des effectifs}.
Aujourd'hui, les transferts de tâches produisent un effet nul par construction : ils sont classés \enquote{pas dans le modèle} (§\ref{pro:leviers}).

\paragraph{La difficulté de fond : besoin non observé, recours observé.}
Le besoin de soins n'est pas observé. Ce qu'on observe en données ouvertes, c'est le \emph{recours} : des actes remboursés par spécialité d'exécutant et par tranche d'âge du bénéficiaire (Open DAMIR, Assurance maladie, cité par la dérivation).
Or le recours est contraint par l'offre : là où les professionnels sont rares, on consomme moins, faute d'accès, et pas forcément faute de besoin.
Prendre le recours observé pour le besoin rend l'indice de couverture en partie circulaire : il sous-estime le besoin là où l'offre est faible, précisément là où la question se pose.
La dérivation mesure $r$ en 2026 et le suppose constant ensuite ; elle ne traite pas ce biais. Les ansätze ci-dessous l'ordonnent du moins cher au plus fidèle.

\paragraph{Ansätze.}
M0 et M1 sont ceux de la dérivation ; M2 à M4 sont proposés par le présent rapport comme programme de recherche, sans prix établi.
\begin{description}
\item[M0, proportionnel aux habitants.] $B\propto\mathrm{pop}(T,t)$. C'est déjà la densité pour 100\,000\mesure{} habitants de la carte : l'indice offre/besoin vaut densité/densité\textsubscript{2026}. Coût nul. Il ne voit ni le vieillissement, ni les transferts. Il reste le cas limite de tous les autres.
\item[M1, recours standardisé par l'âge.] $B(p,T,t)=\sum_{a}\mathrm{pop}(T,t,a)\,r(p,a)$, avec $r(p,a)$ le recours national par habitant et par tranche d'âge en 2026 (Open DAMIR), supposé constant (la dérive du recours, L31, est exclue). Indice de couverture
$C(p,T,t)=\dfrac{\mathrm{offre}(p,T,t)\times\mathrm{activit\acute{e}}/B(p,T,t)}{\mathrm{offre}(p,T,2026)/B(p,T,2026)}$, qui vaut 1\derive{} en 2026. $C<1$ se lit \enquote{couverture moindre qu'en 2026}, dans le lexique neutre.
\item[M2, recours non contraint.] Estimer $r$ là où l'offre ne contraint pas (par exemple dans les territoires les mieux dotés, ou par une frontière de recours en fonction de la densité), puis l'appliquer partout. Il corrige la circularité de M1, au prix d'une hypothèse forte : que le recours des territoires bien dotés mesure le besoin, et non une demande induite par l'offre.
\item[M3, épidémiologique.] Besoin $=\sum_{\text{pathologie}}$ prévalence $\times$ actes requis par cas. Indépendant de l'offre par construction ; demande des prévalences territoriales et des normes de prise en charge, dont aucune n'est chargée dans le dépôt.
\item[M4, spatial (famille 2SFCA).] Accessibilité à deux étapes à aire flottante : l'offre de chaque commune est rapportée à la population pondérée par l'âge de son aire de patientèle, puis sommée sur l'aire de recours de chaque commune, avec une décroissance par la distance. C'est la famille de l'indicateur APL de la DREES et de l'IRDES (exploration \texttt{v2/\allowbreak{}docs/\allowbreak{}lab/\allowbreak{}archive-\allowbreak{}v1/\allowbreak{}explorations/\allowbreak{}drees-\allowbreak{}projections-\allowbreak{}and-\allowbreak{}apl.md} ; méthode résumée dans~\citep{e-bahroun2025}). Il résout la frontière départementale, que M0 à M3 ignorent, mais il demande le grain communal, que la carte n'affiche pas (§\ref{pro:archi}).
\end{description}

\paragraph{Transferts de tâches.}
Dans tous les ansätze, un transfert $\tau$ déplace une part du besoin d'une profession A vers une profession B à partir de la date de décision :
\[
B'(A)=B(A)\,(1-s\,a),\qquad B'(B)=B(B)+s\,a\,B(A)\,k ,
\]
où $a$ est la part du besoin de A concernée par l'acte transféré, $k$ convertit le temps de l'acte de A vers B, et $s$ est le levier. $a$ et $k$ doivent être sourcés ou étiquetés \enquote{supposé}.
Les effectifs projetés ne changent pas : le jumeau ne modélise pas d'installation qui répondrait au besoin, et N9 a mesuré un effet nul de la densité sur les installations (\texttt{decisions.md}, axe 3).
Trois leviers sont visés : vaccination en officine (L23), infirmiers en pratique avancée (L25), accès direct au kiné (L27).

\paragraph{Critères d'acceptation et test d'abandon (dérivation).}
\begin{itemize}
\item Test d'abandon (environ 30\suppose{}~min) : comparer M1 et M0 en 2050, par département, pour les IDE et les MK. Si l'écart reste sous 2\suppose{}\,\% dans plus de 90\suppose{}\,\% des départements, la pondération par l'âge est abandonnée et le module se réduit aux transferts.
\item Avec $r$ constant, $C$ égale densité/densité\textsubscript{2026} à $10^{-9}$\suppose{} près ; avec tous les $\tau$ nuls, effectifs et densités identiques au commit \texttt{23ed102}.
\item Conservation : $\sum B$ avant et après un transfert diffère exactement de $s\,a\,B(A)\,(k-1)$, à $10^{-6}$\suppose{} près.
\item Trois leviers de transfert branchés de bout en bout ; recalcul de la carte sous 150\suppose{}~ms au 95\textsuperscript{e} centile ($n\geq30$\suppose{}) ; moins de 5\suppose{}~ms ajoutées au recalcul.
\item \enquote{Accès direct au kiné à Clermont-Ferrand en 2035} donne un effectif inchangé et un indice de couverture modifié, et l'écran le dit.
\item Ce qui tuerait le plan : le test d'abandon, ou l'impossibilité de sourcer $a$ et $k$ pour aucun des trois transferts. Les leviers resteraient alors \enquote{pas dans le modèle}, avec la raison.
\end{itemize}

\paragraph{L'APL DREES comme cible externe.}
L'APL est publiée à la commune, millésimes 2022 à 2024, sous Licence Ouverte, pour les infirmiers, les kinés, les sages-femmes et les dentistes ; il n'y a pas d'APL des pharmaciens ni des spécialistes (exploration v1, métadonnées du jeu lues, fichiers non téléchargés).
Reproduire l'APL 2024 avec M4, sur les professions communes, donnerait une cible externe de validation, analogue à ce que DD76~\citep{e-dd76} et ER1319~\citep{e-er1319} sont pour les effectifs.
Deux obstacles sont relevés : l'APL borne l'activité de chaque professionnel à partir de données d'activité (SNIIR-AM, DCIR) dont l'équipe ne dispose pas ; et la DREES avertit que l'APL ne compare pas deux millésimes en tenant compte du vieillissement, ce qui en fait une cible de niveau, pas de tendance.
La formule exacte (fonction de décroissance, bornes) est dans les archives \enquote{Méthode concertée}, non ouvertes.

\paragraph{Question de périmètre : la médecine générale côté source.}
La médecine générale est exclue du sujet, et la dérivation l'exclut du module.
Or les trois transferts visés (vaccination, infirmiers en pratique avancée, accès direct au kiné) touchent, selon toute vraisemblance, des actes aujourd'hui assurés ou prescrits en médecine générale : c'est une hypothèse du présent rapport, que le dépôt ne source pas.
La dérivation admet les spécialistes du seul côté \enquote{source}, à titre de contexte. La même règle pour la médecine générale (son besoin $B(A)$ comme source d'un transfert, sans projection de son offre) rendrait L23, L25 et L27 calculables ; elle étendrait aussi le périmètre au-delà du sujet.
C'est une décision de l'équipe, non prise (tableau~\ref{ouv:tab}).
\endgroup
